\documentclass[11pt,a4paper]{article}
\usepackage{amsmath,amssymb,amsthm}
\usepackage{geometry}
\usepackage{tikz}
\usepackage{tcolorbox}
\usepackage{enumitem}
\usepackage{xcolor}
\usepackage{hyperref}
\usepackage{booktabs}
\usepackage{array}
\usepackage{listings}

% Page setup
\geometry{margin=1in}
\setlength{\parindent}{0pt}
\setlength{\parskip}{0.5em}

% Define colors
\definecolor{theoremblue}{RGB}{0,51,102}
\definecolor{definitiongreen}{RGB}{0,102,51}
\definecolor{exampleorange}{RGB}{204,102,0}
\definecolor{keyconceptpurple}{RGB}{102,0,102}
\definecolor{proofbrown}{RGB}{139,69,19}
\definecolor{questionred}{RGB}{180,0,0}
\definecolor{codegray}{RGB}{245,245,245}

% Define theorem environments
\tcbuselibrary{theorems,skins,breakable,listings}

\newtcbtheorem[number within=section]{definition}{Definition}
{colback=green!5,colframe=definitiongreen,fonttitle=\bfseries,breakable}{def}

\newtcbtheorem[number within=section]{theorem}{Theorem}
{colback=blue!5,colframe=theoremblue,fonttitle=\bfseries,breakable}{thm}

\newtcbtheorem[number within=section]{keyconcept}{Key Concept}
{colback=purple!5,colframe=keyconceptpurple,fonttitle=\bfseries,breakable}{key}

\newtcbtheorem[number within=section]{questionbox}{Question and Answer}
{colback=red!5,colframe=questionred,fonttitle=\bfseries,breakable}{qa}

% Custom commands
\newcommand{\R}{\mathbb{R}}
\newcommand{\Z}{\mathbb{Z}}
\newcommand{\E}{\mathbb{E}}

\title{\textbf{Sustainable Life Management Model}\\
\large Questions and Answers on Model Formulation and Solution Approach\\
\large February 6, 2026}
\author{Technical Documentation}
\date{}

\begin{document}

\maketitle
\tableofcontents
\newpage

%==============================================================================
\section{Solution Methodology}
%==============================================================================

\begin{questionbox}{Is the Problem Trivially Solvable by Greedy Algorithm?}{greedy}
\textbf{Question:} If we want to maximize expected discounted utility return, does solving the problem become trivial given that resources are allocated to tasks ordered by their weights? Can this be solved by a greedy algorithm? Is this the current setting and formulation?

\textbf{Answer:} \textbf{No}, the problem is \textbf{not} trivially solvable by a greedy algorithm, despite the objective function being linear in the allocation variables $h_{it}$.

\textbf{Why a greedy approach seems plausible:}

In a simple single-period linear knapsack with one constraint, a greedy algorithm that orders items by value-to-weight ratio and packs them in decreasing order achieves optimality for the LP relaxation. Since the current objective is linear in $h_{it}$, one might compute a per-unit return for each task and allocate greedily.

\textbf{Why greedy fails in practice:}

\begin{enumerate}
    \item \textbf{State-dependent coupling across periods.} The priority weight $w_I(t, g_t) = w_I^0 \cdot (1 - g_t) \cdot \exp(-\rho_I t/T)$ depends on $g_t$, which accumulates from $\sum_{\tau=1}^{t} h_{1,\tau}$. Allocating more hours to Task 1 today reduces its future weight through $(1 - g_t)$. A greedy algorithm cannot anticipate this intertemporal feedback.
    
    \item \textbf{Multi-dimensional knapsack structure.} Each period has three simultaneous constraints: time ($H^{\max}$), energy ($E^{\max}$), and emotional capacity ($M^{\max}$). Multi-dimensional knapsack is NP-hard; greedy does not guarantee optimality even in the single-period case.
    
    \item \textbf{Minimum allocation requirements.} The constraint $h_i^{\min} \cdot x_{it} \leq h_{it}$ means a task cannot receive a fractional amount below its minimum. This creates discrete jumps in feasible allocations that preclude continuous greedy solutions.
    
    \item \textbf{Stochastic recovery coupling.} Recovery constraints link emotional load to available time through $r_{ts} \geq \beta_s \cdot z_{ts}$. Choosing high-emotion tasks reduces time available for other tasks in a scenario-dependent manner. Greedy ranking ignores this coupling.
\end{enumerate}

\textbf{The current formulation:}

The current formulation is a multi-period MILP with state-dependent weights, multiple resource constraints, and stochastic recovery. This places it in the class of NP-hard combinatorial optimization problems. The MILP formulation solved by Gurobi's branch-and-bound is the appropriate approach, not a greedy heuristic.
\end{questionbox}

\begin{questionbox}{Is the Model Solved via Dynamic Programming?}{dp}
\textbf{Question:} Is the model solved in a dynamic programming way, i.e., forward then backward? What does the implementation codebase in \texttt{/src/} show?

\textbf{Answer:} \textbf{No}, the model is \textbf{not} solved using dynamic programming (forward-backward induction). Instead, it is formulated and solved as a \textbf{monolithic mixed-integer linear program (MILP)} using the Gurobi commercial solver.

\textbf{Evidence from the codebase:}
\begin{enumerate}
    \item \textbf{Optimizer class} (\texttt{src/model/optimizer.py}): The model is built as a single Gurobi \texttt{Model} object that contains all decision variables for all time periods simultaneously:
    \begin{itemize}
        \item All variables $x_{it}$, $h_{it}$, $r_{ts}$, $z_{ts}$ for $t = 1, \ldots, T$ are created at once
        \item All constraints across all periods are added to one model
        \item The objective function spans the entire horizon
    \end{itemize}
    
    \item \textbf{Solver class} (\texttt{src/model/solver.py}): The \texttt{Solve()} method simply calls \texttt{self.Model.optimize()}, which invokes Gurobi's branch-and-bound algorithm on the complete formulation.
    
    \item \textbf{No stage decomposition}: There is no evidence of:
    \begin{itemize}
        \item Value function computation
        \item Backward induction over stages
        \item State space enumeration
        \item Bellman equation evaluation
    \end{itemize}
\end{enumerate}

\textbf{Should it be solved via DP?} The problem has characteristics that could theoretically support a DP approach (temporal structure, state variables $g_t$), but the current MILP formulation offers several advantages:
\begin{itemize}
    \item \textbf{Tractability}: With $O(|\mathcal{I}| \cdot T \cdot |\mathcal{S}|)$ variables, the problem is small enough for direct MILP solution (achieves optimality within 60 seconds).
    \item \textbf{State space}: DP would require discretizing the continuous state space $(G_t, \sigma_t)$, introducing approximation errors.
    \item \textbf{Implementation simplicity}: MILP leverages mature commercial solvers without custom algorithm development.
\end{itemize}
\end{questionbox}

%==============================================================================
\section{Uncertainty Modeling}
%==============================================================================

\begin{questionbox}{How is Stochastic Uncertainty Considered?}{uncertainty}
\textbf{Question:} How is stochastic/uncertainty considered in the model beyond the definition/setting? How does it impact the solution process?

\textbf{Answer:} The model employs a \textbf{scenario-based stochastic programming approach} to handle uncertainty. The uncertainty is operationalized through the following mechanisms:

\textbf{1. Scenario Set Definition:}
\begin{itemize}
    \item Discrete scenario set $\mathcal{S} = \{1, 2, 3, 4\}$
    \item Each scenario $s$ specifies recovery hours $\beta_s \in \{1, 2, 3, 4\}$
    \item Probabilities $p_s = \{0.4, 0.3, 0.2, 0.1\}$
\end{itemize}

\textbf{2. Impact on Objective Function:}

The objective computes the \textbf{expected value} over scenarios:
$$\max \sum_{s \in \mathcal{S}} p_s \left[ \sum_{t=1}^T \delta^t \sum_{i \in \mathcal{I}} U_{it}(h_{it}, g_t) \right]$$

From \texttt{src/model/objective.py}:
\begin{verbatim}
for Scen in range(config.SCENARIOS):
    Prob = config.RECOVERY_PROBS[Scen]
    Objective += Prob * (0.95 ** Day) * DayUtility
\end{verbatim}

\textbf{3. Impact on Constraints:}

Scenario-indexed constraints ensure feasibility across all scenarios:
\begin{itemize}
    \item \textbf{Time availability}: $\sum_{i \in \mathcal{I}} h_{it} + r_{ts} \leq H^{\max}$ for all $t, s$
    \item \textbf{Recovery triggering}: $r_{ts} \geq \beta_s \cdot z_{ts}$ for all $t, s$
\end{itemize}

This creates \textbf{scenario-dependent recovery requirements} while the primary allocation decisions $h_{it}$ remain first-stage (here-and-now) variables that do not depend on which scenario realizes.

\textbf{4. Solution Implications:}
\begin{itemize}
    \item Decisions must be robust to all recovery scenarios
    \item Worst-case scenario ($\beta_4 = 4$ hours) constrains maximum productive time
    \item Expected value objective balances performance across scenarios
\end{itemize}
\end{questionbox}

\begin{questionbox}{Where is Uncertainty Explicitly Modeled?}{whereuncertain}
\textbf{Question:} Where in the model does it explicitly consider uncertainty and how is it modeled?

\textbf{Answer:} Uncertainty appears explicitly in three model components:

\textbf{1. Stochastic Recovery Constraints} (Lines 200-202 in Formulation.tex):
\begin{align}
\sum_{i \in \mathcal{I}} m_i \cdot h_{it} &\geq 0.8 \cdot M^{\max} - M \cdot (1 - z_{ts}) \quad \forall t \in \mathcal{T}, s \in \mathcal{S} \\
r_{ts} &\geq \beta_s \cdot z_{ts} \quad \forall t \in \mathcal{T}, s \in \mathcal{S}
\end{align}

\textbf{Interpretation:}
\begin{itemize}
    \item First constraint: If emotional load exceeds 80\% of capacity, the binary indicator $z_{ts}$ is forced to 1
    \item Second constraint: When $z_{ts} = 1$, recovery time $r_{ts}$ must be at least $\beta_s$ hours
    \item The big-M formulation models the logical implication: high emotional load $\Rightarrow$ recovery required
\end{itemize}

\textbf{2. Time Availability with Scenario-Indexed Recovery}:
$$\sum_{i \in \mathcal{I}} h_{it} + r_{ts} \leq H^{\max} \quad \forall t \in \mathcal{T}, s \in \mathcal{S}$$

This ensures that productive time plus scenario-dependent recovery does not exceed available hours.

\textbf{3. Expected Value Objective}:
$$\E_s \left[ \sum_{t=1}^T \delta^t \sum_{i \in \mathcal{I}} U_{it}(h_{it}, g_t) \right] = \sum_{s \in \mathcal{S}} p_s \cdot \text{Utility under scenario } s$$

\textbf{Model Classification:}

This is a \textbf{two-stage stochastic program with recourse}:
\begin{itemize}
    \item \textbf{First-stage decisions}: Task selections $x_{it}$ and allocations $h_{it}$ (made before uncertainty resolves)
    \item \textbf{Second-stage/recourse decisions}: Recovery hours $r_{ts}$ (depend on scenario)
    \item \textbf{Relatively complete recourse}: Any first-stage decision can be made feasible by adjusting recovery
\end{itemize}
\end{questionbox}

\begin{questionbox}{Should Stochastic Returns Be Considered in Addition to Stochastic Size?}{stochreturns}
\textbf{Question:} For the uncertainty part, should stochastic returns also be considered in addition to stochastic size/requirement (1, 2, 3, 4 hours)? Is this already considered given the decay function, or not?

\textbf{Answer:} \textbf{No}, the current model does \textbf{not} model stochastic returns. The decay function captures a different phenomenon. Stochastic returns should be added for a more complete uncertainty treatment.

\textbf{What the decay function captures (deterministic, not stochastic):}

The weight functions
\begin{align}
w_I(t, g_t) &= w_I^0 \cdot (1 - g_t) \cdot \exp(-\rho_I t/T) \\
w_Y(t) &= w_Y^0 \cdot (1 + \rho_Y t/T)
\end{align}
model \textit{systematic, deterministic} trends in priority over time and state. They do not model \textit{random fluctuations} in how productive a task hour is. The decay function is known in advance and identical across all scenarios.

\textbf{What stochastic returns would capture (random, scenario-dependent):}

In reality, the return from one hour of work varies due to:
\begin{itemize}
    \item Cognitive state fluctuations (alertness, focus quality)
    \item Environmental interruptions (unexpected meetings, distractions)
    \item Task complexity variation (some days a problem is harder than expected)
    \item Learning curve effects (breakthrough insights vs. dead ends)
\end{itemize}

\textbf{Extended formulation with stochastic returns:}

Let $\tilde{\alpha}_{is}^d$ denote the return rate for task $i$ in dimension $d \in \{I, Y, K, O\}$ under scenario $s$. The extended utility becomes:
\begin{align}
\tilde{U}_{its}(h_{it}, g_t) = &\; w_I(t, g_t) \cdot \tilde{\alpha}_{is}^I \cdot h_{it} + w_Y(t) \cdot \tilde{\alpha}_{is}^Y \cdot h_{it} \nonumber \\
&+ w_K(t) \cdot \tilde{\alpha}_{is}^K \cdot h_{it} + w_O(t) \cdot \tilde{\alpha}_{is}^O \cdot h_{it} \nonumber \\
&- \lambda_1 \cdot e_i \cdot h_{it} - \lambda_2(g_t) \cdot m_i \cdot h_{it}
\end{align}

\textbf{Relationship between decay and stochastic returns:}

\begin{center}
\begin{tabular}{lll}
\toprule
Aspect & Decay Function & Stochastic Returns \\
\midrule
Nature & Deterministic & Random \\
Varies with & Time $t$ and state $g_t$ & Scenario $s$ \\
Models & Priority evolution & Productivity fluctuation \\
Known when? & Before optimization & Revealed after decision \\
\bottomrule
\end{tabular}
\end{center}

The two are \textit{orthogonal}: decay governs how much we \textit{want} to do a task, while stochastic returns govern how much \textit{value} we get from doing it. With both stochastic returns and stochastic recovery sizes, the problem becomes a stochastic knapsack with uncertainty in both the ``return'' and ``size'' dimensions.
\end{questionbox}

%==============================================================================
\section{Multi-Agent Extension and Advanced Algorithms}
%==============================================================================

\begin{questionbox}{Should Multiple Agents Be Considered?}{multiagent}
\textbf{Question:} The problem setting is designed for one employee/agent. Should multiple employees/agents be considered to make the problem more realistic and applicable? Does this become a multiple knapsack with stochastic return and size? Is it impactful for the introduction/background?

\textbf{Answer:} \textbf{Yes}, extending to multiple agents significantly enhances practical applicability and transforms the problem into a \textbf{multiple knapsack problem with stochastic returns and sizes}.

\textbf{Why multiple agents matter:}

\begin{enumerate}
    \item \textbf{Organizational realism.} Real workplaces involve teams where multiple employees share resources (meeting rooms, supervisory time, collaborative task capacity). Individual optimization without coordination leads to resource contention and suboptimal organizational outcomes.
    
    \item \textbf{Coupling constraints.} Shared resources create coupling across agents:
    $$\sum_{j \in \mathcal{J}} h_{ijt} \leq C_i^{\max} \quad \forall i \in \mathcal{I}_{\text{shared}}, t \in \mathcal{T}$$
    These constraints prevent the multi-agent problem from decomposing into independent single-agent subproblems.
    
    \item \textbf{Problem classification.} With multiple agents, each agent represents a ``knapsack'' with individual capacity constraints (time, energy, emotional load), and tasks represent ``items'' that must be packed. With stochastic returns $\tilde{\alpha}_{is}^d$ and stochastic sizes $\beta_s$, this is precisely a \textit{multiple knapsack problem with stochastic returns and sizes}.
\end{enumerate}

\textbf{Multi-agent formulation:}

Let $\mathcal{J} = \{1, \ldots, |\mathcal{J}|\}$ be the set of agents. The extended objective is:
$$\max \quad \sum_{j \in \mathcal{J}} \sum_{s \in \mathcal{S}} p_s \left[ \sum_{t=1}^T \delta^t \sum_{i \in \mathcal{I}} \tilde{U}_{ijts}(h_{ijt}, g_{jt}) \right]$$

Subject to agent-specific constraints:
\begin{align}
\sum_{i \in \mathcal{I}} h_{ijt} + r_{jts} &\leq H_j^{\max} \quad \forall j, t, s \\
\sum_{i \in \mathcal{I}} e_i \cdot h_{ijt} &\leq E_j^{\max} \quad \forall j, t \\
\sum_{i \in \mathcal{I}} m_i \cdot h_{ijt} &\leq M_j^{\max} \quad \forall j, t
\end{align}

\textbf{Impact on introduction/background:}

\textbf{Yes}, the multi-agent extension is impactful for the introduction. The existing literature review already discusses organizational burnout, workforce scheduling, and team productivity. The multi-agent extension provides:
\begin{itemize}
    \item A direct bridge from individual well-being research to organizational resource management
    \item Connection to the JD-R and COR theories at the team level
    \item Relevance to the gig economy and algorithmic management literature
    \item A richer problem class (multiple knapsack) that justifies advanced algorithmic investigation
\end{itemize}

\textbf{Problem size scaling:}

For $|\mathcal{J}|$ agents, the problem size grows to $O(|\mathcal{J}| \cdot |\mathcal{I}| \cdot T \cdot |\mathcal{S}|)$ variables. With $|\mathcal{J}| = 10$, $|\mathcal{I}| = 7$, $T = 30$, $|\mathcal{S}| = 4$: approximately 8,400 binary variables and 16,800 constraints. This tenfold increase motivates decomposition algorithms.
\end{questionbox}

\begin{questionbox}{Can Advanced Algorithms Beat Direct Gurobi Solution?}{algorithms}
\textbf{Question:} Would the multiple knapsack with stochasticity be solved using more advanced solution algorithms such as branch-and-cut, column generation, and Benders decomposition, and then beat the performance or solution found by Gurobi?

\textbf{Answer:} \textbf{Potentially yes}, particularly for the multi-agent extension. The problem structure creates natural decomposition opportunities that specialized algorithms can exploit.

\textbf{1. Benders Decomposition:}

\textit{Why it fits:} The two-stage stochastic structure separates naturally:
\begin{itemize}
    \item \textbf{Master problem}: First-stage allocation decisions $x_{ijt}, h_{ijt}$ (integer)
    \item \textbf{Subproblems}: For each scenario $s$, determine recovery $r_{jts}$ (continuous LP)
\end{itemize}

Optimality cuts from subproblems iteratively tighten the master. When $|\mathcal{S}|$ is large (many scenarios), Benders avoids solving the full deterministic equivalent and can be significantly faster.

\textit{Expected advantage:} Strong when the number of scenarios grows beyond what direct MILP can handle efficiently (e.g., $|\mathcal{S}| > 50$ with continuous distributions discretized finely).

\textbf{2. Column Generation (Dantzig-Wolfe Decomposition):}

\textit{Why it fits:} The multi-agent formulation has a \textbf{block-diagonal structure}:
\begin{itemize}
    \item Each agent's constraints form an independent block
    \item Coupling constraints (shared resources) link agents at the master level
    \item Each ``column'' represents a complete schedule for one agent
\end{itemize}

The master problem selects columns (agent schedules) satisfying coupling constraints. The pricing subproblem generates improving schedules for individual agents, each of which is a tractable single-agent MILP.

\textit{Expected advantage:} Strong when $|\mathcal{J}|$ is large. The LP relaxation bound from column generation is at least as tight as the direct LP relaxation, often significantly tighter.

\textbf{3. Branch-and-Cut:}

\textit{Why it fits:} The knapsack substructure admits well-known valid inequalities:
\begin{itemize}
    \item \textbf{Cover inequalities}: For each knapsack constraint, identify minimal covers and add cutting planes
    \item \textbf{Lifted cover inequalities}: Strengthen cover cuts by coefficient lifting
    \item \textbf{Gomory mixed-integer cuts}: Derive from fractional LP relaxation solutions
    \item \textbf{Surrogate relaxation}: Combine multiple knapsack constraints into a single surrogate constraint for bound computation
\end{itemize}

\textit{Expected advantage:} Tighter LP relaxations reduce the branch-and-bound tree size. Most effective when integrated with Benders or column generation.

\textbf{When will they beat Gurobi?}

\begin{center}
\begin{tabular}{lll}
\toprule
Method & Beats Gurobi when & Unlikely to beat when \\
\midrule
Benders & $|\mathcal{S}| \gg 4$ & Small scenario count \\
Column gen. & $|\mathcal{J}| \gg 1$ & Single agent \\
Branch-cut & Weak LP relaxation & Gurobi cuts suffice \\
\bottomrule
\end{tabular}
\end{center}

\textbf{Important caveat:} Gurobi internally implements many of these techniques (branch-and-cut with Gomory cuts, cover cuts, etc.). The advantage of custom implementations lies in exploiting \textit{problem-specific} structure that generic solvers cannot recognize automatically, such as the agent-decomposable block structure or the two-stage scenario structure. For the current single-agent formulation with $|\mathcal{S}| = 4$, direct Gurobi solution in 60 seconds is likely near-optimal. The decomposition advantage emerges at scale.
\end{questionbox}

%==============================================================================
\section{Key Concepts from Integer Programming}
%==============================================================================

\begin{keyconcept}{Mixed-Integer Linear Programming (MILP)}{milp}
A MILP has the general form:
\begin{align*}
\max \quad & c^T x + d^T y \\
\text{s.t.} \quad & Ax + By \leq b \\
& x \in \Z^p, \quad y \in \R^q
\end{align*}

where $x$ are integer variables and $y$ are continuous variables.

\textbf{Solution approaches:}
\begin{itemize}
    \item \textbf{Branch-and-bound}: Systematically explores the solution space using LP relaxations
    \item \textbf{Branch-and-cut}: Adds cutting planes to tighten LP relaxations
    \item \textbf{Dynamic programming}: Stage-wise optimization with value functions (not used here)
\end{itemize}

The SLM model uses MILP because:
\begin{enumerate}
    \item Binary selection variables $x_{it} \in \{0,1\}$ for task activation
    \item Continuous allocation variables $h_{it} \in \R_+$ for hours
    \item Linear objective and constraints enable efficient solution
\end{enumerate}
\end{keyconcept}

\begin{keyconcept}{Stochastic Programming}{stochprog}
Stochastic programming optimizes decisions under uncertainty:

\textbf{Two-stage formulation:}
\begin{align*}
\min_{x} \quad & c^T x + \E_\xi[Q(x, \xi)] \\
\text{s.t.} \quad & Ax = b, \quad x \geq 0
\end{align*}

where $Q(x, \xi)$ is the optimal value of the second-stage problem given first-stage decision $x$ and scenario $\xi$.

\textbf{Scenario-based approximation:}

Replace expectation with weighted sum over discrete scenarios:
$$\E_\xi[Q(x, \xi)] \approx \sum_{s \in \mathcal{S}} p_s \cdot Q(x, s)$$

This converts the stochastic program to a deterministic equivalent problem (DEP), which is exactly what the SLM formulation implements.
\end{keyconcept}

\begin{keyconcept}{Dynamic Programming vs. MILP}{dpvsmilp}
\textbf{Dynamic Programming:}
\begin{itemize}
    \item Decomposes problem by stages/time periods
    \item Computes value functions via backward induction
    \item Optimal policy obtained by forward simulation
    \item Suffers from ``curse of dimensionality'' with continuous states
\end{itemize}

\textbf{Bellman equation:}
$$V_t(s) = \max_{a \in A(s)} \left\{ r(s, a) + \gamma \sum_{s'} P(s'|s,a) V_{t+1}(s') \right\}$$

\textbf{MILP approach:}
\begin{itemize}
    \item Formulates entire horizon as single optimization
    \item Exploits problem structure via branch-and-bound
    \item Handles continuous variables naturally
    \item Leverages commercial solver technology
\end{itemize}

\textbf{When to use each:}
\begin{itemize}
    \item \textbf{DP}: Large discrete state spaces, Markov structure, need for policy
    \item \textbf{MILP}: Moderate problem size, mixed integer-continuous variables, commercial solver access
\end{itemize}

The SLM problem falls into the MILP-favorable category due to continuous state variables and tractable problem size.
\end{keyconcept}

%==============================================================================
\section{Summary}
%==============================================================================

\begin{keyconcept}{Key Takeaways}{summary}
\begin{enumerate}
    \item \textbf{Solution method}: The SLM model is solved as a monolithic MILP using Gurobi's branch-and-bound algorithm, \textbf{not} via dynamic programming or greedy algorithms.
    
    \item \textbf{Non-triviality}: Despite linear per-period returns, the problem is NP-hard due to state-dependent weight coupling, multi-dimensional knapsack constraints, minimum allocation requirements, and stochastic recovery coupling. Greedy ranking by weight does not yield optimal solutions.
    
    \item \textbf{Uncertainty handling}: Scenario-based stochastic programming with discrete scenarios for both recovery requirements (stochastic size) and task returns (stochastic returns). The objective maximizes expected utility, and constraints ensure feasibility across all scenarios.
    
    \item \textbf{Decay vs. stochastic returns}: The decay function $w_I(t, g_t)$ captures deterministic temporal priority evolution. Stochastic returns $\tilde{\alpha}_{is}^d$ capture random productivity fluctuations. These are orthogonal and both should be modeled.
    
    \item \textbf{Explicit uncertainty locations}:
    \begin{itemize}
        \item Scenario-indexed recovery variables $r_{ts}$
        \item Stochastic return coefficients $\tilde{\alpha}_{is}^d$
        \item Binary high-load indicators $z_{ts}$
        \item Time availability constraints
        \item Expected value objective
    \end{itemize}
    
    \item \textbf{Multi-agent extension}: Extending to $|\mathcal{J}|$ agents yields a multiple knapsack problem with stochastic returns and sizes, with $O(|\mathcal{J}| \cdot |\mathcal{I}| \cdot T \cdot |\mathcal{S}|)$ variables. Coupling constraints through shared resources prevent independent decomposition.
    
    \item \textbf{Advanced algorithms}: Benders decomposition (scenario structure), column generation (agent structure), and branch-and-cut (knapsack polytope) can exploit problem structure. These are expected to outperform direct Gurobi solution as team size and scenario count grow.
    
    \item \textbf{Two-stage structure}: First-stage (here-and-now) decisions for task allocation; second-stage (recourse) decisions for recovery adjustment.
    
    \item \textbf{Tractability}: Single-agent problem with $O(|\mathcal{I}| \cdot T \cdot |\mathcal{S}|)$ is tractable for direct MILP. Multi-agent extension motivates decomposition for scalability.
\end{enumerate}
\end{keyconcept}

\end{document}
